\begin{abstract}
Compound sparsity---simultaneously pruning attention heads and evicting KV cache entries---is essential for efficient LLM inference but produces super-additive quality degradation whose structure is poorly understood.
We find that compound sparsity severity is not a fixed constant but an emergent property of the interaction among eviction strategy, reasoning phase, and sparsity level.
On DeepSeek-R1-Distill-Qwen-7B (MATH-500), importance-based eviction exhibits a monotonically increasing phase--severity gradient: the paired median difference $\Delta$ rises from $0.39$ at sparsity $s{=}0.15$ to $8.48$ at $s{=}0.30$ (Spearman $\rho{=}1.0$; all Bonferroni-corrected $p < 0.005$).
Random eviction develops a detectable gradient only above $s \approx 0.20$, at $2.7$--$17{\times}$ smaller magnitude.
A scoring-independent control (recency eviction) produces no positive phase gradient: $\Delta$ vanishes at $s \leq 0.25$ and remains non-significant at $s{=}0.30$, confirming that the interaction is attention-score-specific rather than an inherent property of token removal.
The pattern replicates on Qwen2.5-7B-Instruct ($\Delta$: $0.07 \to 3.16$; $\rho{=}1.0$), shows consistent phase gradient direction on GSM8K, and generalizes to LlamaForCausalLM ($n{=}250$; $d_z{=}0.81$--$1.15$, all $p < 10^{-29}$).
Per-phase Gini analysis identifies \emph{importance discrimination collapse}---the degeneration of head importance scores toward uniformity under compound stress---as a correlational diagnostic of this severity divergence.
Prior work characterizes compound degradation as phase-agnostic; we show that phase modulates severity by $15$--$17{\times}$ on DeepSeek and $2.7$--$3.5{\times}$ on Qwen, specifically driven by attention-score-based selection. These findings establish phase-gated analysis as a principled diagnostic framework with direct implications for phase-aware budget allocation.
\end{abstract}
